%%%%%%%%%%  scan idx 84  |  folio 71  %%%%%%%%%%
\SGAPrintPage{71}%
\sgaStableTarget{sga6.e00.bibliographical-references}\subsection*{Bibliographical references:}
\begin{enumerate}[label={[\arabic*]}]
\item \emph{Hirzebruch}: Neue topologische Methoden in der algebraischen Geometrie.
\item \emph{Serre}: Faisceaux algébriques cohérents, Ann. of Maths.
\item \emph{Grothendieck}: Sur les faisceaux algébriques cohérents, in Séminaire Cartan 1956/57.
\end{enumerate}

\sgaStableTarget{sga6.e00.proof-of-the-riemann-roch-theorem}\section*{Proof of the Riemann--Roch theorem}
\begin{center}
(Grothendieck, 1 November 1957)
\end{center}

\sgaStableTarget{sga6.e00.reminder-of-notation}\subsection*{* Reminder of notation}
\SGAAuthorFootText{unmarked}{The passages between asterisks \textsuperscript{*}\dots\textsubscript{*} are interpolations by the copyist.}
[These are those of the ``report on Riemann--Roch'', cited as R.R.R. in what follows.]

All varieties considered are nonsingular and quasi-projective (i.e.\ isomorphic to a locally closed subvariety of a projective space). The base field is algebraically closed, of arbitrary characteristic.

For such a variety \(X\), we denote by \(K(X)\) the abelian group generated by the classes of coherent sheaves on \(X\) (modulo the identification of an extension with a sum); the alternating sum of Tor makes \(K(X)\) a ring. If \(\underline F\) is a sheaf (resp.\ a vector bundle \(E\), resp.\ a sub%
%%%%%%%%%%  scan idx 85  |  folio 72  %%%%%%%%%%
\SGAPrintPage{72}%
variety \(Y\) of \(X\)), we denote by \(\gamma_X(\underline F)\) (resp.\ \(\gamma_X(E)\), resp.\ \(\gamma_X(Y)\)) the element of \(K(X)\) corresponding to \(\underline F\) (resp.\ to \(\underline E\), resp.\ to \(\underline{O}_Y\)). (In fact, we often write \(\underline F\), \(E\) instead of \(\gamma_X(\underline F)\)\dots).

We denote by \(A(X)\) the graded ring of classes of cycles on \(X\) (for rational equivalence---cf.\ Chow); \(A^i(X)\) consists of cycles of codimension \(i\).

For a vector bundle \(E\), we denote by \(C(E)=\sum C^i(E)\), \(C^i(E)\in A^i(X)\), its Chern class (defined as the inverse image of Schubert cycles); by linearity we deduce the \(C^i(y)\), for every \(y\in K(X)\). If \(C(y)=\prod(1+a_i)\), we put, following Hirzebruch:
\[
  \ch(y)=\sum e^{a_i},\qquad
  T(y)=\prod\frac{a_i}{1-e^{-a_i}}.
\]
These are elements of \(A(X)\otimes\mathbb Q\).

When \(y\) is the tangent bundle of \(X\), we write \(T(X)\) instead of \(T(y)\).

For a vector bundle \(G\), we denote by \(\lambda^pG\) its \(p\)-th exterior power, and put \(\lambda_{-1}(G)=\sum(-1)^p\lambda^p(G)\); this is an element of \(K(X)\).

For a proper regular map \(f\) from a variety \(Y\) to a variety \(X\), we define (additive) homomorphisms:
\[
  f^*:A(X)\longrightarrow A(Y)\qquad\text{and}\qquad
  f_*:A(Y)\longrightarrow A(X).
\]

The first is homogeneous of degree \(0\) and multiplicative. The second is homogeneous of degree \(\dim.X-\dim.Y\).

We also define homomorphisms:
\[
  f^!:K(X)\longrightarrow K(Y)\qquad\text{and}\qquad
  f_!:K(Y)\longrightarrow K(X).
\]

For a coherent sheaf \(\underline F\) on \(X\), \(f^!(\underline F)\) equals, by definition, the alternating sum of the \(\Tor_p^{\underline{O}_X}(\underline{O}_Y,\underline F)\). For vector bundles, this is simply the inverse-image operation.

%%%%%%%%%%  scan idx 86  |  folio 73  %%%%%%%%%%
\SGAPrintPage{73}%
For a coherent sheaf \(G\) on \(Y\), \(f_*(\underline G)\) is coherent on \(X\). Denoting by \(R^pf_*\) the \(p\)-th derived functor of the functor \(f_*\) thus defined, we have, by definition, \(f_!(\underline G)=\sum(-1)^p R^pf_*(\underline G)\).\quad\textsubscript{*}

\sgaStableTarget{sga6.e00.riemann-roch-theorem}\textbf{RIEMANN--ROCH THEOREM.} \emph{Let \(f:Y\longrightarrow X\) be a proper morphism} (\textsuperscript{*} regular map \textsubscript{*}) \emph{from a nonsingular quasi-projective variety \(Y\) to another \(X\), and let \(y\in K(Y)\). We have:}
\sgaStableTarget{sga6.e00.eq.1}\begin{equation}\tag{1}
  f_*\bigl(\ch(y)T(Y)\bigr)=\ch(f_!(y))T(X)
  \quad\text{\emph{in} } A(X)\otimes\mathbb Q.
\end{equation}

\emph{Proof} (\textsuperscript{*} abridged \textsubscript{*}).

\sgaStableTarget{sga6.e00.lem.1}\textbf{LEMMA 1.} \emph{Consider proper morphisms of varieties \(Z\xrightarrow{g}Y\xrightarrow{f}X\). If the theorem \(\RR\) is true for \(f\) and \(g\), it is true for \(f\circ g\). If it is true for \(fg\) and \(g\), it is true for \(f\). If \(f_*\) is injective, and the theorem \(\RR\) is true for \(fg\) and \(f\), it is true for \(g\).}

Trivial. (\textsuperscript{*} In fact, when we assume \(\RR\) true for \(fg\) and \(g\), we must also assume that \(g_!\) is surjective to be able to assert that \(\RR\) is true for \(f\). \textsubscript{*})

As a proper morphism is the composite of an injection and a projection \(X\times P\longrightarrow X\) (where \(P\) is a projective space), we are reduced to treating these two cases separately. The second, as we have seen, reduces to a standard calculation, by proposition 2.13 (``cellular decomposition'') of R.R.R. We shall therefore assume \(Y\subset X\), with the injection morphism \(i\). Formula (1) then becomes:
\sgaStableTarget{sga6.e00.eq.2}\begin{equation}\tag{2}
  \ch(i_!(y))=i_*\bigl(\ch(y)T(E)^{-1}\bigr),
\end{equation}
where \(E=T_{X/Y}\) is the normal bundle of \(Y\) in \(X\).

%%%%%%%%%%  scan idx 87  |  folio 74  %%%%%%%%%%
\SGAPrintPage{74}%
\sgaStableTarget{sga6.e00.lem.2}\textbf{LEMMA 2.} \emph{Assume that there exists on \(X\) a vector bundle \(G\) of rank \(p=\codim_XY\), such that:}
\sgaStableTarget{sga6.e00.eq.3}\begin{equation}\tag{3}
  \gamma_X(Y)=i_!(1)=\lambda_{-1}(G)\,;\qquad
  i^!(G)=\check E\,;\qquad
  (-1)^pC^p(G)=Y.
\end{equation}
\emph{If we then have:}
\sgaStableTarget{sga6.e00.eq.4}\begin{equation}\tag{4}
  y=i^!(x),\qquad \text{\emph{with} }x\in K(X),
\end{equation}
\emph{formula (2) is valid.}

The proof is immediate from the fact that \(\ch\) is a ring homomorphism, and that we have:
\sgaStableTarget{sga6.e00.eq.5}\begin{equation}\tag{5}
  \ch(\lambda_{-1}(G))=(-1)^pC^p(G)T(\check G)^{-1}.
\end{equation}

\sgaStableTarget{sga6.e00.cor.unnamed.occ6}\textbf{COROLLARY.} \emph{\(\RR\) is true if \(Y\) is a hypersurface and \(y\) is of the form \(i^!(x)\), with \(x\in K(X)\).}

\sgaStableTarget{sga6.e00.lem.3}\textbf{LEMMA 3.} \emph{Assume that \(X=Y\times P\), with \(P\) a projective space, and that \(i\) is the injection \(u\longrightarrow (u,v_0)\) \(v_0\in P\). Then \(i_*\) is injective and \(\RR\) is true.}

This is pure standard calculation.

Assume again that \(X,Y,y\) are arbitrary (still with \(Y\subset X\)). Let \(X'\) be the variety obtained by blowing up \(X\) along \(Y\), and \(Y'\) the inverse-image hypersurface of \(Y\). We have a diagram of canonical maps
\[
\begin{tikzcd}
Y \arrow[r,"i"] & X\\
Y' \arrow[u,"g"] \arrow[r,"j"] & X' \arrow[u,"f"']
\end{tikzcd}
\]
(\(j\) being the canonical map). The inverse image of \(E\) on \(Y'\) has a subbundle of rank \(1\), whose dual is denoted by \(L\). Moreover, \(L^{-1}\) is also the restriction to \(Y'\) of the bundle \(L(Y')\) defined on \(X'\) by the divisor \(Y'\). We put \(F=E/L^{-1}\). Elementary calculations give:

%%%%%%%%%%  scan idx 88  |  folio 75  %%%%%%%%%%
\SGAPrintPage{75}%
\sgaStableTarget{sga6.e00.lem.4}\textbf{LEMMA 4.} \emph{We have the formulas (where \(p=\codim_XY\)):}
\sgaStableTarget{sga6.e00.eq.6}\sgaStableTarget{sga6.e00.eq.7}\sgaStableTarget{sga6.e00.eq.8}\sgaStableTarget{sga6.e00.eq.9}\begin{align}
  f^!i_!(y)&=j_!\bigl(\lambda_{-1}(\check F)y\bigr). \tag{6}\\
  \lambda_{-1}(\check F)&\equiv \gamma^{p-1}(E)\lambda^p(\check E)
    \quad\text{mod.}(1-L)\text{\SGAAuthorFootMark{*}}. \tag{7}\\
  f_*(1)&=1,\quad\text{\emph{whence} }f_*f^*=1. \tag{8}\\
  f^*i_*(\alpha)&\equiv j_*\bigl(\alpha C^{p-1}(F)\bigr)
    \quad\text{mod.}\Ker f_*\quad\text{\emph{for every} }\alpha\in A(Y). \tag{9}
\end{align}
\SGAAuthorFootText{*}{This unnecessarily complicated formula can be replaced, for example, by VII 4.1.}

[To simplify notation, we have made \(K(Y)\) operate on \(K(Y')\) and \(A(Y)\) on \(A(Y')\) by means of the homomorphisms \(g^!\) and \(g^*\). {*}In formula (7), it seems that \(\gamma^{p-1}(E)\) denotes \(\lambda^{p-1}(E-2)\), that is, \(\sum_{i+j=p-1}\binom{-p}{j}\lambda^i(E)\); the congruence holds in \(K(Y')\).\textsubscript{*}]

Formula (6) follows from the equality:
\sgaStableTarget{sga6.e00.eq.6bis}\begin{equation}\tag{6 bis}
  \Tor_i^{\underline{O}_X(v)}\bigl(\underline{O}_Y(v),\underline{O}_{Y'}(v')\bigr)
  =
  \wedge^i\check F(v')
  \quad (v'\in Y',\ v=g(v')).
\end{equation}

Formula (7) follows from the formulas
\sgaStableTarget{sga6.e00.eq.10}\begin{equation}\tag{10}
  g_!(L^{-i})=0\quad\text{for }0<i\leq p-1,\qquad g_!(1)=1,
\end{equation}
which themselves follow from K\"unneth and the cohomology of projective spaces.

Formula (8) is trivial; finally, (9) follows from \(g_*(C^{p-1}(F))=1\), which is deduced, for example, from the immediate formulas:
\sgaStableTarget{sga6.e00.eq.10bis}\begin{equation}\tag{10 bis}
  g_*(\xi^i)=0\quad\text{if }0\leq i<p-1,\qquad
  g_*(\xi^{p-1})=1,
\end{equation}
with \(\xi=C^1(L)\in A^1(Y')\).

[In fact, (9) is very probably an identity, and not merely a congruence.]


